% # 第四章 基于形式化验证的处理器微架构安全性证明
\chapter{基于形式化验证的处理器微架构安全性证明}
\label{chp:Formal}

第三章介绍的模糊测试技术能够在硅后阶段通过动态探测有效地检测与挖掘微架构漏洞，但其本质是基于概率的搜索过程，无法对处理器的安全属性提供完备性证明。而在处理器设计的早期阶段，芯片尚以寄存器传输级源码的形式存在，研究者能够获取完整的电路逻辑，从而有条件利用数学手段更加充分地验证处理器的安全属性。因此，形式化验证作为一种基于严谨数学穷举的证明方法，被广泛应用于处理器硅前阶段的安全性证明中。本章聚焦于硅前阶段的处理器微架构安全性证明问题，与第三章的硅后动态探测共同构成本文全生命周期安全性证明体系的两大支柱。

% ============================================================================
% 4.1 引言
% ============================================================================
\section{引言}
\label{sec:ch4_introduction}

在现代处理器设计中，处理器核心不是孤立运行的计算单元，而是通常与缓存、内存控制器、中断控制器、外部设备等多种平台组件集成在一起，构成片上系统（System on Chip, SoC）。如图~\ref{fig:ch4_soc}所示，在这一体系中，处理器核心负责指令执行与控制逻辑，而平台组件则提供存储访问、中断响应、外设通信等基础服务，它们之间通过多种总线协议（如 AMBA AXI、TileLink）进行连接，整体上构成了一个完整的计算设备。

\begin{figure}[htb!]
    \centering
    \includegraphics[width=0.7\textwidth]{figures/content/4.1 片上系统架构.pdf}
    \caption{片上系统（SoC）架构示意图。处理器核心与外设、加速器、内存等平台组件通过总线协议互连，共同构成完整的计算平台。}
    \label{fig:ch4_soc}
\end{figure}

目前，以 ShadowLogic 为代表的形式化验证方案已经能够在 RTL 级别对处理器核心的安全性进行完备证明，在多款 RISC-V 处理器核心上取得了有效验证。然而，处理器核心的安全性不代表整个片上系统的安全性，处理器核心与平台组件之间的交互可能引入新的时序泄露。这类时序泄露源于处理器核心与平台之间的交互过程，处于现有研究方法的验证范围之外。现有形式化验证技术的验证范围主要局限于处理器核心本身，在验证过程中仅针对处理器核心中已有的部分进行安全性验证，忽略了整个处理器设计在后期搭载更多平台组件时，仍然可能引入新的时序泄露。

平台组件（如缓存、中断控制器等）的响应行为往往依赖于运行时状态，可能引入处理器核心隔离验证中不可见的时序通道，导致通过验证的处理器核心在集成到真实平台后仍可能被利用。而将处理器核心与平台合并进行组合验证，又面临平台 RTL 不可得与状态空间爆炸的双重障碍。第\ref{sec:ch4_problem}节将对上述局限性进行详细分析。

\textbf{本文贡献。}针对上述问题，本章提出基于片上系统时序合约的解耦验证框架，包括：

（1）本文提出了片上系统时序合约（SoC Timing Contract）的概念，并设计了相应的形式化定义方法，通过片上系统时序合约分别建模合约下的处理器核心特征与平台特征，系统化描述处理器核心与平台之间的时序关系，设计了针对理想内存、缓存、物理内存保护（Physical Memory Protection, PMP）三类典型平台场景适用的相应时序合约，并基于格理论建立合约之间的强弱偏序关系，为分层验证提供理论依据。

（2）基于片上系统时序合约，本文设计了时序合约插桩方法（Timing Contract Instrumentation），并构建了基于时序合约插桩的双向解耦验证策略。时序合约插桩方法根据不同的片上系统时序合约条例，对待验证单元（Design Under Test, DUT）构造的等价性检查电路（Miter Circuit, Miter 电路）进行插桩，构造不同的平台时序行为，实现在没有接入实际平台的条件下检测处理器核心的安全特性。在此基础上，本文将时序合约插桩进一步分解为处理器时序合约插桩（CPU Timing Contract Instrumenting, CPU\_TCI）与平台时序合约插桩（Platform Timing Contract Instrumenting, Platform\_TCI）两类，分别对处理器核心和平台组件进行独立验证，只要处理器和平台分别通过同一层级的 CPU\_TCI 和 Platform\_TCI 验证，即可保证二者集成后仍然满足该合约规定的安全性，从而实现无需组合验证的安全性保证。

\textbf{本章结构。}第\ref{sec:ch4_problem}节介绍基于软件硬件合约的形式化验证框架及其局限性；第\ref{sec:ch4_framework}节详细介绍基于片上系统时序合约的解耦验证框架；第\ref{sec:ch4_experiment}节进行实验评估；第\ref{sec:ch4_cex_analysis}节对实验中发现的典型反例进行深入分析；第\ref{sec:ch4_summary}节总结本章。


% ============================================================================
% 4.2 问题描述
% ============================================================================
\section{基于软件硬件合约的形式化验证框架}
\label{sec:ch4_problem}

针对第\ref{sec:ch4_introduction}节指出的处理器核心与平台交互中的安全性验证问题，本节介绍现有形式化验证方案的技术路线与验证边界，并分析其局限性。

\subsection{基于乘积电路的合约验证方法}
\label{subsec:ch4_miter}

如第二章所述，软件硬件安全合约通过合约轨迹与硬件轨迹定义程序的安全属性。第三章的模糊测试将合约实例化为动态差分测试判据，而形式化验证则将同一合约实例化为电路级的等价性断言。其中，形式化验证将合约轨迹映射为针对电路属性的前提假设（assume），通过对电路的初始状态进行限定，过滤不满足合约前提的目标程序；将硬件轨迹映射为针对电路属性的安全断言（assert），通过检验电路在不同输入条件下的安全属性是否会被违反，实现对电路整体安全性的断言。如图~\ref{fig:ch4_miter}所示，基于上述思路，现有的基于乘积电路的合约验证方法可分为乘积电路层、合约层与验证层三个部分。

\textbf{乘积电路层。}软件硬件安全合约的安全属性本质上是一个二元关系，需要比较两组不同秘密输入下的可观测行为是否一致。然而，形式化验证所依赖的模型检测器只能对单一电路实例进行推理，无法直接跨电路比较多条轨迹。为此，研究者引入乘积电路（Miter Circuit），将两条执行轨迹实例化在同一个电路系统中，使模型检测器能够通过单一电路的推理间接完成两条轨迹的比较。具体而言，乘积电路设计一个顶层电路模块，同时实例化同一处理器设计两份。以 RISC-V 处理器为例，与安全属性相关的数据输入主要来源于内存的初始内容：程序 $P$ 存储在指令内存中，数据存储在数据内存中，其中公开数据记为 $M_{pub}$，秘密数据记为 $M_{sec}$。乘积电路将程序 $P$ 与公开数据 $M_{pub}$ 作为两个副本的共享输入，同时为两个副本各自初始化不同的秘密数据 $M_{sec}$ 与 $M_{sec}'$ 作为独立输入。在形式化验证中，程序 $P$ 与数据内存的初始值通常被设为符号化（symbolic）的抽象值，使得模型检测器能够自动搜索所有可能的程序与数据组合，而非依赖于人工构造的测试用例。通过这种构造，两个处理器副本使用相同的设计、共享相同的程序与公开输入，电路特性保持一致；而独立的秘密输入又为差分比较提供了条件，使得后续能够在单一电路中探索不同秘密输入下的处理器状态差异。

\textbf{合约层。}合约层是一组附加在乘积电路之上的监视逻辑，通过附加电路实现针对电路特性的约束过滤、观测提取与断言比较，构成具体的安全性证明逻辑。约束过滤对应合约轨迹的验证逻辑，用于过滤合约轨迹不一致的输入组合。若两个副本的架构级行为（如提交地址序列、分支跳转方向）存在差异，则说明程序本身的控制流或访存行为依赖了秘密数据。这种泄露属于软件层面的泄露，并不是硬件漏洞。例如，若程序根据秘密值进行了不同方向的分支跳转，那么即使在一个没有任何微架构侧信道的理想单周期处理器上，攻击者仅通过观察程序的控制流就能推断秘密。这类泄露是软件自身的缺陷，而非硬件漏洞，因此应将其排除在验证范围之外。约束过滤确保验证仅关注在架构级行为一致的前提下，微架构实现是否引入了额外的时序差异，即真正由硬件导致的泄露。观测提取对应硬件轨迹的属性提取逻辑，为后续的安全属性断言做铺垫，从两个副本中提取微架构级的可观测特征。不同的合约层级关注不同的泄露维度，因而定义了不同的可观测特征集合，典型的可观测特征包括指令提交时间序列与访存地址序列等。安全性断言对应硬件轨迹的验证逻辑，正式执行安全属性的验证，比较两个副本提取出的硬件特征是否满足安全属性，进而推断处理器设计的安全性。一个周全的处理器设计需要保证，对于任意输入组合，只要合约轨迹一致（即程序的架构级行为不依赖于秘密数据），其硬件轨迹也必须一致（即微架构级的可观测行为也不泄露秘密数据）。

\textbf{验证层。}整个框架的验证过程通过形式化验证引擎进行穷举搜索。验证引擎将电路特性转化为数学约束，在符号化的输入空间上穷举探索所有可能的输入组合，尝试证明在任意情况下安全属性是否依然能够得到满足。常用的验证引擎包括抽象模型引擎（Abstract Model, AM），适用于推导无界安全性证明；以及有界模型检测引擎（Bounded Model Checking, BMC），适用于在有限步长内搜索反例。验证结果有三种可能：若安全属性断言被验证引擎证明，在所有可达状态下均成立，则处理器满足合约（Proven），其在该合约下的安全性得到证明；若验证引擎找到一组使断言不成立的输入赋值，则返回一个具体的违反反例（CEX），可直接用于攻击场景的分析；若在给定的时间或搜索步长内未能得出确定性结论，则报告验证超时（Timeout）。

\begin{figure}[htb!]
    \centering
    \includegraphics[width=0.85\textwidth]{figures/content/4. 基于乘积电路的合约检测方法.pdf}
    \caption{基于乘积电路的合约验证方法框架图。框架分为乘积电路层（两份处理器副本，共享公开输入、秘密输入独立）、合约层（约束过滤、观测提取与断言比较）和验证层（穷举搜索，输出 Proven/Timeout/CEX）三个部分。}
    \label{fig:ch4_miter}
\end{figure}

\subsection{ShadowLogic 验证框架}
\label{subsec:ch4_baseline}

基于第\ref{subsec:ch4_miter}节介绍的通用验证方法，ShadowLogic 是当前该方向中最具代表性的工作。ShadowLogic 采用常数时间合约（CT）作为安全属性，通过引入影子逻辑（Shadow Logic）实现合约层功能，使得乘积电路验证方法能够在多款真实乱序处理器上落地。

影子逻辑是一组附加在乘积电路两个处理器副本之上的监视电路，不修改处理器内部设计，仅通过层次化信号访问监控处理器的提交阶段，由处理器设计者根据微架构特点编写。在 CT 合约下，影子逻辑实现合约层的三个步骤如下：（1）在约束过滤中，影子逻辑在每条指令提交时比较两个副本的 ISA 架构级行为，包括提交的访存地址、分支跳转方向以及下一条指令的 PC 选择，若上述信号存在差异，则通过 \texttt{invalid\_program} 信号将当前输入标记为无效，确保验证仅覆盖满足 CT 合约前提的程序。（2）在观测提取中，影子逻辑从两个副本中提取两类微架构级可观测特征：指令提交时序（两个副本是否在同一周期提交同一条指令）与访存地址序列（两个副本的 load/store 地址是否一致）。（3）在安全性断言中，影子逻辑检验上述可观测特征是否存在差异，若在约束过滤确认程序满足 CT 合约前提的条件下，两个副本的提交时序或访存地址仍然存在不一致，则判定处理器违反了 CT 安全属性。

在乱序处理器上，两个副本因微架构状态不同（如重排序缓冲区占用、功能单元竞争等），可能在不同的时钟周期提交同一条指令。为实现准确的逐指令比较，ShadowLogic 通过门控时钟（clock gating）进行同步对齐：当一个副本先提交指令而另一个尚未提交时，影子逻辑暂停先提交副本的时钟信号，等待另一个副本追上后再恢复，从而实现逐指令而非逐周期的对齐比较。此外，影子逻辑采用两阶段验证机制以确保指令覆盖的完整性：第一阶段逐周期运行两个副本并监控微架构观测信号，一旦检测到偏差则进入第二阶段；第二阶段等待两个副本流水线中的所有在途指令排空，确保偏差发生时流水线中的每条指令都已被约束过滤检查，排空完成后再判定该偏差是否构成真实的安全属性违反。

ShadowLogic 已在多款 RISC-V 处理器上完成验证，涵盖顺序流水线处理器（Sodor）、简单乱序处理器（SimpleOoO）、超标量乱序处理器（RideCore）以及复杂乱序处理器（BOOM），能够在安全的处理器设计上给出形式化证明（Proven），在不安全的设计上找到具体攻击序列（CEX）。然而，ShadowLogic 的验证范围仅限于处理器核心本身——验证过程中处理器的内存接口被封装在验证单元内部（如 Sodor 处理器的 SodorInternalTile），平台响应行为被隐含假设为固定延迟的理想组件。这一假设在处理器核心集成到实际片上系统后不再成立。

\subsection{局限性分析}
\label{subsec:ch4_limitations}

ShadowLogic 实现了处理器核心的高效安全性证明，但其验证范围与实际部署场景之间存在差距。本节从隔离验证的安全性盲区与组合验证的可扩展性瓶颈两个方面分析现有方案的局限性。

\textbf{局限性一：隔离验证存在安全性盲区。}

现有方案的验证对象是封装了内存的处理器模块（如 Sodor 处理器的 SodorInternalTile），模块内部的内存响应时序固定（1 周期），两个副本的内存行为完全一致。验证通过意味着在这种理想化内存假设下，处理器核心不存在违反安全合约的时序泄露。

然而，处理器核心在实际片上系统中并非孤立运行，它通过总线接口与缓存、中断控制器等平台组件交互。这些平台组件的响应行为不再是固定延迟的——缓存存在地址依赖的命中与缺失时序差异，中断控制器的响应可能受写入数据影响。以缓存为例，乱序处理器的推测执行可能发出秘密依赖的访存地址，在理想内存下不产生时序差异（延迟固定），但在有缓存的平台上，不同地址导致命中或缺失，产生时序差异，进而影响后续指令的提交时间，构成时序侧信道。这条攻击路径在隔离验证中完全不可见。

隔离验证的结论是有条件的——它隐含假设了平台的理想化行为。当这一假设不成立时，通过验证的处理器核心在实际部署中仍可能被攻击。现有方案缺乏对平台交互引入的时序通道的建模能力。

\textbf{局限性二：组合验证面临可扩展性瓶颈。}

将处理器核心与平台组件合并为整体设计并纳入形式化验证范围，理论上可以消除局限性一中的安全性盲区。然而，这种组合验证方法在工程实践中面临三重障碍。

其一，平台 RTL 不可得。在处理器设计的早期阶段，目标平台的完整 RTL 设计通常尚未完成，平台规格往往以接口协议或行为规范的形式存在，无法直接作为验证输入使用。

其二，状态空间爆炸。即使平台 RTL 可得，组合设计的状态空间远大于处理器核心本身，形式化验证的复杂度随电路规模急剧增长。以 ShadowLogic 的实际数据为参考，仅 BOOM 处理器核心本身的安全设计已难以在 7 天内完成证明，若再叠加缓存和中断控制器，验证将更加困难。

其三，验证不可复用。组合验证将处理器与特定平台绑定，每次更换平台配置（如更换缓存策略、增加中断控制器），都需要对整个组合设计重新进行完整验证，之前的验证成果无法复用，验证成本与平台配置数量成正比。

综合两方面局限性：隔离验证覆盖不足，组合验证难以实现。这要求一种新的验证方法，既能在不依赖平台 RTL 的前提下覆盖平台交互引入的时序通道，又能避免组合验证的状态空间爆炸。第\ref{sec:ch4_framework}节将介绍本文提出的基于片上系统时序合约的解耦验证框架，通过合约抽象平台行为实现处理器核心与平台的独立验证。


% ============================================================================
% 4.3 基于片上系统时序合约的解耦验证框架
% ============================================================================
\section{基于片上系统时序合约的解耦验证框架}
\label{sec:ch4_framework}

针对第\ref{subsec:ch4_limitations}节分析的两个局限性，本节提出基于片上系统时序合约的解耦验证框架。其中，时序合约的形式化定义与时序合约插桩方法解决局限性一（隔离验证的安全性盲区），双向解耦验证策略解决局限性二（组合验证的可扩展性瓶颈）。第\ref{subsec:ch4_overview}节介绍总体框架；第\ref{subsec:ch4_contract}节介绍时序合约的形式化定义；第\ref{subsec:ch4_decouple}节给出双向解耦的验证策略；第\ref{subsec:ch4_cpu_tci}节介绍处理器侧验证方法；第\ref{subsec:ch4_plat_tci}节介绍平台侧验证方法。

\subsection{总体框架}
\label{subsec:ch4_overview}

本节介绍基于片上系统时序合约的解耦验证方法的总体框架。

传统的基于软件硬件合约的形式化验证框架主要依赖于构造乘积电路，证明待验证处理器设计的安全性。本文沿用这一验证思想，在此基础上引入片上系统时序合约的概念，将验证范围从处理器核心拓展至片上系统整体，并提出解耦方法论分别证明处理器侧与平台侧的安全性。如图~\ref{fig:ch4_framework}所示，第\ref{sec:ch4_introduction}节概述的两项贡献在本节展开为以下四个递进的方面：

\begin{enumerate}
    \item 本文形式化定义了片上系统时序合约，具体分为 C1 理想内存型平台合约、C2 缓存型平台合约与 C3 物理内存保护型（Physical Memory Protection, PMP）平台合约。三种合约对应了常见的三种平台类型，系统化描述了处理器核心与平台之间可能存在的时序泄露关系。
    \item 基于时序合约，本文提出了双向解耦的验证策略，将片上系统的安全性证明拆分为处理器侧与平台侧两步独立验证。只要两侧分别通过同一合约层级的验证，即可保证集成后的安全性，无需进行组合验证。
    \item 本文提出了处理器侧安全性证明方法（CPU\_TCI）。CPU\_TCI 通过在乘积电路上进行时序合约插桩，在无需获取平台 RTL 代码的条件下模拟平台的时序行为特征，检测处理器核心在特定合约假设下的安全性。
    \item 本文还提出了平台侧安全性证明方法（Platform\_TCI）。Platform\_TCI 通过构造平台组件的乘积电路，在无需获取处理器核心 RTL 代码的条件下，检测具体平台的实际行为是否仅产生合约允许范围内的时序差异，验证平台是否满足特定合约层级的时序约束。
\end{enumerate}

\begin{figure}[htb!]
    \centering
    \includegraphics[width=0.95\textwidth]{figures/content/4.3 基于片上系统时序合约的解耦验证框架.pdf}
    \caption{基于片上系统时序合约的解耦验证框架。左侧为处理器侧验证方法（CPU\_TCI），通过时序合约插桩模拟平台行为；右侧为平台侧验证方法（Platform\_TCI），检查平台是否满足合约约束。两侧共享同一合约层级，各自独立验证后通过组合安全性定理保证集成安全性。}
    \label{fig:ch4_framework}
\end{figure}

\subsection{片上系统时序合约的形式化定义}
\label{subsec:ch4_contract}

本节通过分析常见平台设备类型（理想内存、缓存、内存映射外设）的时序信号特征，总结归纳其可能引入的时序泄露模式，提出片上系统时序合约，将这些时序特征抽象为形式化的合约定义，系统化描述处理器核心与平台之间可能存在的时序泄露关系，为后续的解耦验证提供形式化基础。

\textbf{接口信号模型。}片上系统时序合约基于一个通用的简化内存接口模型设计。该模型包括处理器核心的输出信号与输入信号两类：输出信号为处理器核心向平台发出的请求信号，包括内存请求信号（\texttt{req}）、访存地址（\texttt{addr}）、读写类型（\texttt{we}）以及写入数据（\texttt{wdata}）；输入信号为平台返回给处理器核心的响应信号，包括响应有效信号（\texttt{gnt}）、返回数据（\texttt{rdata}）以及中断信号（\texttt{int}，作为可选扩展）。该接口模型与 HWPE 2.0 接口协议的 req/gnt 握手语义一致，遵循 HWPE 2.0 或语义等价接口规范的平台可以直接复用合约逻辑。片上系统时序合约关注的核心问题是：处理器核心的哪些输出信号被允许影响平台返回的哪些输入信号的时序与数据。

\textbf{合约语法。}本文采用 diamond 算子（$\Diamond$）形式化表达输出信号对输入信号的因果影响关系。一条合约条目的形式为：
\begin{equation}
    \texttt{output\_signal} \rightarrow \Diamond\{\texttt{input\_signals}\}
    \label{eq:ch4_diamond}
\end{equation}
其含义为：若两个处理器副本的该输出信号曾经产生过差异，则合约允许其后续接收到的对应输入信号也存在差异。$\Diamond$ 算子强调的是``最终影响''关系——输出信号在任意时刻产生的差异，都可能在后续任意时刻影响对应的输入信号。一个完整的片上系统时序合约由若干条上述条目组成，每条输出信号对应一组被允许影响的输入信号集合。

\textbf{合约实例。}根据平台组件的复杂程度，本文定义了三类片上系统时序合约，其形式化定义如表~\ref{tab:ch4_contracts}所示，各合约的设计依据如下：

\begin{enumerate}
    \item C1 合约（固定延迟内存）。理想内存的响应时序与访存地址无关，所有访问均为固定延迟，因此仅请求信号被允许影响响应时序。C1 是最强的合约层级，对平台行为的约束最为严格。
    \item C2 合约（地址依赖延迟）。缓存的命中与缺失导致响应延迟依赖于访存地址，读写类型也可能影响缓存行为（如写回策略下的额外延迟），因此 C2 在 C1 基础上额外允许地址和读写类型影响响应时序。C2 较 C1 更弱，允许更多的时序差异。
    \item C3 合约（数据依赖中断）。内存映射外设（如中断控制器）的响应可能受写入数据影响，因此 C3 在 C2 基础上引入条件性约束：仅当访存地址落在外设地址范围（\texttt{periph\_range}）内时，写入数据才被允许影响响应时序和中断信号。该地址范围可通过物理内存保护（PMP）机制限定。C3 较 C2 更弱。
\end{enumerate}

\begin{table}[htb]
    \centering
    \caption{片上系统时序合约定义}
    \label{tab:ch4_contracts}
    \begin{tabular}{lll}
        \hline
        \textbf{合约} & \textbf{平台特征} & \textbf{形式化定义} \\
        \hline
        \multirow{4}{*}{C1} & \multirow{4}{*}{\shortstack[l]{固定延迟内存}} & \texttt{req} $\rightarrow \Diamond$\{\texttt{gnt, rdata}\} \\
        & & \texttt{addr} $\rightarrow \Diamond$\{\texttt{rdata}\} \\
        & & \texttt{we} $\rightarrow \Diamond$\{\texttt{rdata}\} \\
        & & \texttt{wdata} $\rightarrow \Diamond$\{\texttt{rdata}\} \\
        \hline
        \multirow{4}{*}{C2} & \multirow{4}{*}{\shortstack[l]{地址依赖延迟\\（缓存）}} & \texttt{req} $\rightarrow \Diamond$\{\texttt{gnt, rdata}\} \\
        & & \texttt{addr} $\rightarrow \Diamond$\{\texttt{gnt, rdata}\} \\
        & & \texttt{we} $\rightarrow \Diamond$\{\texttt{gnt, rdata}\} \\
        & & \texttt{wdata} $\rightarrow \Diamond$\{\texttt{rdata}\} \\
        \hline
        \multirow{4}{*}{C3} & \multirow{4}{*}{\shortstack[l]{数据依赖中断\\（PMP 外设）}} & \texttt{req} $\rightarrow \Diamond$\{\texttt{gnt, rdata, int}\} \\
        & & \texttt{addr} $\rightarrow \Diamond$\{\texttt{gnt, rdata, int}\} \\
        & & \texttt{we} $\rightarrow \Diamond$\{\texttt{gnt, rdata, int}\} \\
        & & \texttt{wdata} $\rightarrow \Diamond$\{\texttt{gnt, rdata, int}\} \\
        & & \quad (if \texttt{addr} $\in$ \texttt{periph\_range}) \\
        \hline
    \end{tabular}
\end{table}

\textbf{基于格理论的合约偏序关系。}不同的片上系统时序合约之间存在强弱偏序关系。合约 A 强于合约 B，当且仅当合约 A 对平台行为的约束比合约 B 更严格——即合约 A 允许平台产生的时序差异是合约 B 允许范围的子集。上述定义构成格（lattice）结构：C1 最强，要求平台响应时序不依赖于地址，约束最为严格；C2 较弱，放松了地址依赖约束，允许缓存带来的时序差异；C3 更弱，进一步允许写入数据在特定条件下影响响应时序与中断。这一偏序关系具有重要的单调性：若处理器核心在较弱的合约（如 C2，允许更多时序差异）下通过安全性验证，则在更强的合约（如 C1，允许更少时序差异）下也必然安全。验证者可以根据目标平台的实际特征选择对应层级的合约，而无需逐一验证所有更强的合约层级。

\subsection{基于时序合约的解耦验证策略}
\label{subsec:ch4_decouple}

基于第\ref{subsec:ch4_contract}节定义的片上系统时序合约，本节将时序合约作为处理器核心与平台之间的接口规范，构建双向解耦的验证策略。在该策略下，处理器侧验证``处理器核心在合约描述的平台行为下是否安全''，平台侧验证``平台的实际行为是否在合约允许的范围内''。二者通过共享同一合约层级独立验证，无需获取对方的 RTL 设计。为实现这一策略，本文分别针对处理器侧与平台侧构造乘积电路，并在合约层增加时序合约插桩，将安全性证明解耦为两步独立验证。

\textbf{组合安全性定理。}双向解耦验证策略的核心定理为：若处理器核心通过合约 $C_i$ 下的 CPU\_TCI 验证（即处理器核心在满足 $C_i$ 的平台上安全），且平台通过 $C_i$ 下的 Platform\_TCI 验证（即平台满足 $C_i$ 的时序约束），则处理器核心与该平台集成后的片上系统满足安全属性。形式化表述为：
\begin{equation}
    \textit{CPU\_TCI}(C_i) \wedge \textit{Platform\_TCI}(C_i) \implies \textit{Secure}(\text{CPU} \oplus \text{Platform}, C_i)
    \label{eq:ch4_compositional}
\end{equation}
其中 $\textit{CPU\_TCI}(C_i)$ 表示处理器核心通过合约 $C_i$ 下的验证，$\textit{Platform\_TCI}(C_i)$ 表示平台通过合约 $C_i$ 下的验证，等式右侧表示二者集成后的片上系统（$\text{CPU} \oplus \text{Platform}$）在合约 $C_i$ 下是安全的。
该定理使得处理器设计团队和平台设计团队能够各自独立验证，通过共享的合约层级 $C_i$ 作为接口保证集成安全性。具体而言，处理器核心更换防御机制后只需重新执行 CPU\_TCI 验证，无需重新验证平台；平台更换缓存策略后只需重新执行 Platform\_TCI 验证，无需重新验证处理器核心，验证成果可跨设计变更复用。

\textbf{验证流程。}基于上述组合安全性定理，双向解耦验证的具体流程如下：其一，根据目标平台的特征选择对应的合约层级 $C_i$（如带缓存的平台选择 C2）；其二，执行 CPU\_TCI($C_i$) 验证处理器核心在合约 $C_i$ 假设下的安全性；其三，执行 Platform\_TCI($C_i$) 验证平台是否满足合约 $C_i$ 的时序约束；其四，若两方均通过验证，则处理器核心与该平台的集成安全性得到保证，无需进行组合验证。

\textbf{完备性论证。}上述验证流程的安全性保证源于两侧验证的互补性。CPU\_TCI($C_i$) 在乘积电路中通过无约束的符号化平台响应，穷举了合约 $C_i$ 允许范围内的所有可能平台行为，因此其验证结论覆盖了满足 $C_i$ 的所有平台，而非某一特定平台实现。Platform\_TCI($C_i$) 则确认具体平台的实际行为确实处于合约 $C_i$ 的允许范围之内。二者结合，前者保证了处理器核心对合约范围内所有平台行为的安全性（完备性），后者保证了实际平台确实处于该范围之内（合规性），从而排除了因拆分验证而遗漏安全隐患的可能。

\textbf{不兼容分析。}当验证未通过时，通过不同合约层级下 CPU\_TCI 与 Platform\_TCI 结果的交叉比对可以定位问题来源。例如，假设处理器核心的验证结果为 CPU\_TCI($C_1$) 通过但 CPU\_TCI($C_2$) 失败，说明处理器核心仅在满足 C1（固定延迟内存）的平台上安全，无法抵御缓存引入的地址依赖时序差异。同时，假设目标平台的验证结果为 Platform\_TCI($C_1$) 失败但 Platform\_TCI($C_2$) 通过，说明该平台存在地址依赖的时序差异（如缓存命中与缺失），不满足 C1 但满足 C2。交叉比对两侧结果：处理器核心需要 C1 平台，但实际平台仅满足 C2，二者不兼容。基于定位结果，设计者可以选择修复处理器核心（增强防御使其在 C2 下也安全）、修复平台（替换为满足 C1 的固定延迟内存）或同时修复二者，第\ref{sec:ch4_experiment}节将通过具体实验验证上述不兼容分析流程。

\subsection{处理器侧验证方法（CPU\_TCI）}
\label{subsec:ch4_cpu_tci}

针对第\ref{subsec:ch4_decouple}节提出的解耦验证策略中处理器侧的验证需求，本节介绍 CPU\_TCI 的具体实现方法。

\textbf{乘积电路构造。}与 ShadowLogic 使用封装了内部内存的 Tile 模块不同，CPU\_TCI 将验证单元替换为暴露内存接口的处理器核心模块（Core），使处理器核心与平台之间的接口信号直接可见。CPU\_TCI 构造两个 Core 副本，共享相同的指令内存（执行同一程序），独立的数据内存（秘密数据不同，公开数据相同）。两个副本的数据差异来源于独立的秘密输入 $M_{sec}$ 与 $M_{sec}'$，时序差异则由时序合约插桩逻辑控制。

\textbf{插桩核心组件。}在上述乘积电路的两个处理器核心副本之间，CPU\_TCI 附加时序合约插桩逻辑，监控两个副本的输出信号差异并控制副本 2 的平台响应输入。时序合约插桩由三个核心组件构成：

其一，XOR 差异检测。对两个副本的每个处理器核心输出信号进行逐位异或比较，检测两个副本是否曾经产生过不同的输出。对于多位信号，采用归约或（reduction OR）将多位异或结果合并为单一差异指示信号。

其二，sticky-one 算子。实现合约定义中的 $\Diamond$ 算子，其门级电路实现如图~\ref{fig:ch4_sticky_one}所示，由一个 OR 门与一个 D 触发器构成反馈回路。当差异检测信号触发时，OR 门输出为 1 并被 D 触发器锁存，后续即使差异信号消失，触发器的反馈输出仍使 OR 门保持输出 1，从而实现``一旦置位便永久保持''的语义，仅在系统复位时归零。

\begin{figure}[htb!]
    \centering
    \includegraphics[width=0.7\textwidth]{figures/content/4.3 sticky-one.pdf}
    \caption{Sticky-one 算子。左侧为简化符号表示，右侧为等价的门级电路实现。输入 diff\_signal 为 XOR 差异检测的输出，输出 sticky\_signal 一旦被置为 1 便通过反馈回路永久保持，仅在复位（RST）时归零。}
    \label{fig:ch4_sticky_one}
\end{figure}

其三，插桩 MUX 控制。根据合约允许的时序差异关系，决定两个处理器副本是否接收相同的平台响应。副本 1 始终接收共享的平台响应作为基准，副本 2 的响应由插桩 MUX 控制：当合约允许某输入信号存在差异时（即对应的 sticky-one 已被置位），副本 2 接收无约束的（unconstrained）响应；否则接收与副本 1 相同的共享响应。

图~\ref{fig:ch4_cpu_tci_circuit}展示了以 C2 合约为例的处理器侧时序合约插桩的完整信号流。两个处理器核心副本（DUT 1 与 DUT 2）的输出信号经 XOR 差异检测与 sticky-one 记录后，由插桩逻辑根据合约配置生成 allow\_diff 控制信号，最终通过 MUX 阵列决定 DUT 2 接收共享响应或无约束响应。

\begin{figure}[htb!]
    \centering
    \includegraphics[width=0.85\textwidth]{figures/content/4.3 处理器侧时序合约插桩逻辑.pdf}
    \caption{处理器侧时序合约插桩逻辑（以 C2 合约为例）。DUT 1 与 DUT 2 的输出信号（req, we, addr, wdata）经 XOR 差异检测后由 Sticky-one 算子永久记录，插桩逻辑根据合约配置组合 sticky 信号生成 allow\_diff 控制信号，MUX 阵列据此决定 DUT 2 接收共享平台响应或无约束响应。}
    \label{fig:ch4_cpu_tci_circuit}
\end{figure}

\textbf{不同合约的插桩配置。}不同层级的合约对应不同的插桩逻辑配置，每种合约的插桩配置直接反映其描述的平台时序特征，合约定义与电路实现一一对应。表~\ref{tab:ch4_tci_config}汇总了三种合约下各 allow 控制信号的组合逻辑。

\begin{table}[htb]
    \centering
    \caption{不同合约层级下的插桩控制信号配置}
    \label{tab:ch4_tci_config}
    \begin{tabular}{llll}
        \hline
        \textbf{合约} & $\textit{allow\_gnt\_diff}$ & $\textit{allow\_rdata\_diff}$ & $\textit{allow\_int\_diff}$ \\
        \hline
        C1 & $s_{\textit{req}}$ & $s_{\textit{req}}$ & 0 \\
        C2 & $s_{\textit{req}} \vee s_{\textit{addr}} \vee s_{\textit{we}}$ & $s_{\textit{req}} \vee s_{\textit{addr}} \vee s_{\textit{we}} \vee s_{\textit{wdata}}$ & 0 \\
        C3 & C2 $\vee$ $s_{\textit{wdata}}^{*}$ & C2 $\vee$ $s_{\textit{wdata}}^{*}$ & $s_{\textit{wdata}}^{*}$ \\
        \hline
    \end{tabular}
    \begin{flushleft}
    \small 注：$s_x$ 为 $\textit{sticky\_x}$ 的简写；$s_{\textit{wdata}}^{*}$ 表示条件性检测，仅当 $\texttt{addr} \in \texttt{periph\_range}$ 时生效。
    \end{flushleft}
\end{table}

\begin{enumerate}
    \item C1 合约（固定延迟内存）。C1 合约仅允许请求信号影响响应时序（见第\ref{subsec:ch4_contract}节），对应的插桩配置为 $\textit{allow\_gnt\_diff} = s_{\textit{req}}$，即地址、读写类型和写入数据的差异不参与 $\textit{allow\_gnt\_diff}$ 的计算，仅被允许影响返回数据。
    \item C2 合约（地址依赖延迟）。C2 合约额外允许地址和读写类型影响响应时序（见第\ref{subsec:ch4_contract}节），对应的插桩配置为 $\textit{allow\_gnt\_diff} = s_{\textit{req}} \vee s_{\textit{addr}} \vee s_{\textit{we}}$。写入数据仍仅被允许影响返回数据，因此 $\textit{allow\_rdata\_diff}$ 额外包含 $s_{\textit{wdata}}$，而 $\textit{allow\_gnt\_diff}$ 不包含。
    \item C3 合约（数据依赖中断）。C3 合约在 C2 基础上引入条件性约束（见第\ref{subsec:ch4_contract}节），对应的插桩配置为：仅当 $\texttt{addr} \in \texttt{periph\_range}$ 时，$s_{\textit{wdata}}$ 才被纳入 $\textit{allow\_gnt\_diff}$ 的计算，同时新增 $\textit{allow\_int\_diff} = s_{\textit{wdata}}^{*}$ 控制中断信号。此外，C3 还需要在影子逻辑中区分程序发起的 PC 变化（分支跳转）和平台发起的 PC 变化（中断响应），确保中断引起的控制流差异不被误判为程序违反合约。PMP 约束作为额外的软件假设，禁止秘密数据写入外设地址范围。
\end{enumerate}

\subsection{平台侧验证方法（Platform\_TCI）}
\label{subsec:ch4_plat_tci}

针对第\ref{subsec:ch4_decouple}节提出的解耦验证策略中平台侧的验证需求，本节介绍 Platform\_TCI 的具体实现方法。与 CPU\_TCI 从处理器核心视角出发不同，Platform\_TCI 从平台视角出发，验证具体的平台实现是否满足特定合约层级的时序约束。

\textbf{平台乘积电路构造。}Platform\_TCI 构造两个平台副本（不包含处理器核心），两个副本接收处理器核心的输出信号作为输入，产生返回给处理器核心的响应信号作为输出。验证过程检查合约不允许的信息流在平台实现中是否存在。

\textbf{验证逻辑。}Platform\_TCI 的目标是证明平台产生的所有可能时序差异都在合约允许的范围内。具体而言，合约规定了当某些输出信号（output\_signal）不同时，哪些输入信号（input\_signal）被允许不同。Platform\_TCI 需要验证的是反向命题：若某个输出信号不同，则合约\textbf{未列出}的输入信号必须保持一致。若存在合约未允许的输入信号因输出信号差异而产生变化，则说明平台引入了合约范围之外的时序通道，验证失败。

\textbf{验证示例。}以下通过具体示例说明 Platform\_TCI 的验证过程：

\begin{enumerate}
    \item 验证常规缓存是否满足 C2 合约。C2 合约未允许 \texttt{wdata} 影响 \texttt{gnt}，因此需要检验：当两个缓存副本的 \texttt{wdata} 不同时，\texttt{gnt} 是否仍保持一致。构造两个缓存副本，给予相同的 \texttt{req}、\texttt{addr} 和 \texttt{we}，但不同的 \texttt{wdata}。由于常规缓存的响应时序仅依赖于地址（命中或缺失），不依赖于写入数据，\texttt{gnt} 一致，验证通过（Proven）。
    \item 验证同一常规缓存是否满足 C1 合约。C1 合约未允许 \texttt{addr} 影响 \texttt{gnt}，因此需要检验：当两个缓存副本的 \texttt{addr} 不同时，\texttt{gnt} 是否仍保持一致。给予相同的 \texttt{req}，但不同的 \texttt{addr}。由于不同地址导致不同的命中与缺失结果，\texttt{gnt} 不一致，验证失败（CEX）——缓存的地址依赖时序超出了 C1 的允许范围。
    \item 验证中断控制器是否满足 C2 合约。C2 合约未允许 \texttt{wdata} 影响 \texttt{int}，因此需要检验：当两个中断控制器副本的 \texttt{wdata} 不同时，\texttt{int} 是否仍保持一致。给予相同的 \texttt{req} 和 \texttt{addr}，但不同的 \texttt{wdata}。由于不同的写入数据触发不同的中断行为，\texttt{int} 不一致，验证失败（CEX）——中断控制器的数据依赖行为超出了 C2 的允许范围。
\end{enumerate}

上述示例表明，同一平台在不同合约层级下具有不同的合规性。Platform\_TCI 的验证结果与 CPU\_TCI 的结果交叉比对后，即可判断处理器核心与平台的兼容性。


% ============================================================================
% 4.4 实验评估
% ============================================================================
\section{实验评估}
\label{sec:ch4_experiment}

本节对第\ref{sec:ch4_framework}节提出的解耦验证框架进行实验评估。针对第\ref{subsec:ch4_decouple}节提出的``发现问题→定位问题→指导修复''三步闭环验证流程，本节分别从时序合约插桩的有效性（第\ref{subsec:ch4_eval_tci}节）、合约层级的定位能力（第\ref{subsec:ch4_eval_locate}节）、解耦修复的可行性（第\ref{subsec:ch4_eval_fix}节）三个方面进行验证，并在第\ref{subsec:ch4_eval_scale}节将解耦验证与组合验证进行系统对比。

\subsection{实验设置}
\label{subsec:ch4_setup}

\textbf{实验环境。}全部实验在一台配备 Intel Xeon Gold 6258R 处理器（28 核 56 线程，主频 2.70 GHz，最大睿频 4.00 GHz）、39 MB L3 缓存的服务器上运行。由于形式化验证为确定性过程，相同输入必然产生相同的验证结果（Proven/CEX/Undetermined），不存在随机性，验证时间的波动仅来源于系统调度，因此本文报告单次运行的壁钟时间。

\textbf{实验对象——处理器模块。}本文选取两款架构特征互补的开源 RISC-V 处理器进行验证，二者均来自 ShadowLogic 的原始实验，确保与现有工作的可比性：

\begin{enumerate}
    \item SimpleOoO。一款支持 6 条指令（LI、ADD、MUL、LD、ST、BR）的 4 级流水线乱序处理器，采用 4 条目重排序缓冲区，实现了两种推测执行防御策略：NoFwd 策略允许推测状态下的 load 指令发射但在提交前不转发结果，Delay 策略在推测状态下不发射 load 指令，等到指令提交后再执行访存操作。选择 SimpleOoO 用于 C1/C2 合约的验证，因为缓存时序通道的触发依赖于推测执行机制（推测 load 以秘密数据作为地址导致缓存命中/缺失差异），且其内置的多种防御策略便于对比不同防御强度下的验证结果。
    \item Sodor。一款 2 级流水线的顺序处理器，实现 RV32I 指令集，具备完整的中断接口（含中断输入与 CSR 处理逻辑）。选择 Sodor 用于 C3 合约的验证，因为中断时序通道的触发不依赖于推测执行（已提交的 store 指令向中断控制器写入秘密数据即可触发），且 Sodor 作为常用的形式化验证基准处理器，其验证结果具有更广泛的参考价值。
\end{enumerate}

两款处理器形成互补：SimpleOoO 展示了推测执行与缓存交互产生的时序泄露风险，Sodor 展示了即使在顺序执行的处理器上，中断通道同样可能引入新的时序泄露。

为适配本文的解耦验证框架，需要将处理器核心与平台充分解耦，使核心模块仅包含纯粹的计算逻辑，不内嵌任何平台组件。这要求将处理器的验证单元从封装内存的模块替换为暴露内存接口的核心模块，使时序合约插桩能够在核心与平台之间的接口上施加控制。具体而言，SimpleOoO 通过将内部内存移至外部并增加数据内存读写接口实现，流水线、重排序缓冲区及防御机制等核心逻辑保持不变；Sodor 直接使用已有的 Core 模块替代 SodorInternalTile，Core 模块本身不做任何修改。所有新增的时序合约插桩逻辑与影子逻辑均编写在顶层验证模块中，不修改处理器核心的 RTL 源码。

\textbf{实验对象——平台模块。}实验涉及以下三类平台组件：

\begin{enumerate}
    \item 理想内存。为 ShadowLogic 原始实验中使用的固定延迟内存（1 周期），所有访问的响应时序与地址、数据均无关，对应 C1 合约假设。
    \item 缓存模块，包括 Regular Cache 与 Cache-S 两种实现。Regular Cache 为单条目直接映射缓存，包含缓存标签（tag）与命中/缺失判断逻辑，命中延迟 1 周期、缺失延迟 3 周期，其响应时序依赖于访存地址。Cache-S 是 Regular Cache 的安全替代实现，去除了缓存标签与命中/缺失判断逻辑，所有访问均为固定 1 周期延迟，提供与 Regular Cache 相同的接口但消除了地址依赖的时序差异。在真实系统中，安全缓存可通过分区、随机化等技术实现，本文的 Cache-S 采用最简单的固定延迟设计。
    \item 中断控制器。用于 Sodor 处理器的 C3 合约验证，通过内存映射寄存器接收写入数据并产生中断信号。
\end{enumerate}

\textbf{验证工具。}实验使用 Cadence JasperGold 形式化验证平台，采用抽象模型引擎（AM）进行安全性证明。全部实验采用常数时间合约（CT），观测模型为提交的访存地址序列与指令提交时序。在 JasperGold 中，时序合约通过 SystemVerilog Assertions（SVA）转化为具体的验证属性：合约轨迹的约束过滤通过 \texttt{assume} 语句实现，将不满足合约前提的输入组合排除在验证范围之外；安全性断言通过 \texttt{assert} 语句实现，要求两个处理器副本的可观测输出在所有时钟周期上保持一致。处理器的指令内存与数据内存初始值设为符号化抽象值，由验证引擎自动搜索所有可能的输入组合。

% =============================================================================
% 4.4.2 时序合约插桩的有效性验证
% =============================================================================
\subsection{时序合约插桩的有效性验证}
\label{subsec:ch4_eval_tci}

本节验证时序合约插桩方法能否在不获取平台 RTL 的条件下，等效检出平台交互引入的时序泄露。实验以 SimpleOoO 处理器的 NoFwd 防御配置为对象，如图~\ref{fig:ch4_verification_scope}所示设计三组对比实验：（a）隔离验证作为基线，展示现有方案的验证边界；（b）组合验证作为参照，展示平台引入的真实安全风险；（c）TCI 插桩验证作为本文方法，验证其能否在无需平台 RTL 的条件下检出与组合验证一致的泄露。

\begin{figure}[htb!]
    \centering
    \includegraphics[width=0.95\textwidth]{figures/content/4.4 verification_scope.pdf}
    \caption{三种验证方法的验证范围对比。（a）隔离验证将处理器核心与理想内存封装为整体进行验证，不考虑外部平台组件；（b）组合验证将处理器核心与真实缓存合并纳入验证范围；（c）TCI 插桩验证通过时序合约插桩模拟平台时序行为，无需缓存 RTL 即可覆盖平台引入的时序通道。}
    \label{fig:ch4_verification_scope}
\end{figure}

首先验证现有隔离验证方案的局限性。在隔离验证场景下（图~\ref{fig:ch4_verification_scope}(a)），使用 ShadowLogic 对 NoFwd 配置的 SimpleOoO 进行验证，处理器内存被建模为固定延迟的理想组件，验证结果为 Proven（耗时 212 秒），表明在隔离环境下处理器核心不存在违反 CT 合约的时序泄露。然而，当将同一处理器核心与常规缓存组合后进行验证（图~\ref{fig:ch4_verification_scope}(b)），验证结果为 CEX（反例，16 个时钟周期，耗时 60.55 秒），即形式化验证引擎在 16 个周期内找到了一条违反安全合约的执行轨迹。缓存的命中与缺失产生了地址依赖的时序差异，构成了新的时序侧信道。该结果证实了第\ref{subsec:ch4_limitations}节分析的局限性一：隔离验证通过的处理器核心在集成缓存后不再安全。

随后验证 TCI 方法的有效性。使用 CPU\_TCI 在 $C_2$ 合约下对同一 NoFwd 处理器核心进行验证（图~\ref{fig:ch4_verification_scope}(c)），无需接入常规缓存的 RTL 代码，仅通过时序合约插桩模拟缓存的时序行为特征。验证结果同样为 CEX（16 个时钟周期，耗时 32.91 秒），与组合验证检出了相同的泄露，且验证速度更快。表~\ref{tab:ch4_tci_validity}汇总了上述三组实验的结果。

\begin{table}[htb]
    \centering
    \caption{时序合约插桩有效性验证结果。实验按论证逻辑排列：隔离验证（基线）→ 组合验证（参照）→ TCI 插桩验证（本文方法）。}
    \label{tab:ch4_tci_validity}
    \begin{tabular}{cllrl}
        \hline
        \textbf{序号} & \textbf{平台/合约} & \textbf{方法} & \textbf{时间} & \textbf{结果} \\
        \hline
        1 & 隔离（理想内存） & ShadowLogic & 212s & Proven \\
        2 & 常规缓存 & 组合验证 & 60.55s & CEX (16 cycles) \\
        3 & CPU\_TCI($C_2$) & TCI 插桩 & 32.91s & CEX (16 cycles) \\
        \hline
    \end{tabular}
\end{table}

上述三组实验形成递进式论证：隔离验证的 Proven 结果表明现有方案能够证明处理器核心在理想内存假设下的安全性，但该证明不涵盖外部平台组件。组合验证的 CEX 结果证实了隔离验证通过的设计在集成缓存后重新暴露时序泄露风险。TCI 插桩验证的 CEX 结果表明本文方法在不获取缓存 RTL 的条件下，检出了与组合验证一致的泄露（相同的 16 周期反例），且验证时间从 60.55 秒缩短至 32.91 秒（提速约 1.8 倍），这是因为 TCI 使用符号化的无约束响应替代了真实缓存电路，避免了缓存状态对验证状态空间的扩展。

此外，对比隔离验证与 TCI 插桩验证的验证时间可以评估 TCI 插桩引入的额外开销。隔离验证耗时 212 秒，而 TCI 插桩验证耗时 32.91 秒，TCI 插桩不仅未增加验证时间，反而因为使用去除了内部内存的 Core 模块，缩小了验证单元的电路规模，使验证开销降低至隔离验证的约 $1/6$。这表明时序合约插桩方法在拓展验证覆盖范围的同时，不引入显著的性能损失，其额外开销可忽略不计。

% =============================================================================
% 4.4.3 合约层级定位能力验证
% =============================================================================
\subsection{合约层级定位能力验证}
\label{subsec:ch4_eval_locate}

本节验证时序合约的层级体系能否精准定位处理器核心与平台之间的不兼容来源。实验分别在处理器侧和平台侧进行多合约层级的 TCI 验证，通过交叉比对两侧结果判断问题出在处理器侧还是平台侧。本节依次展示 CPU\_TCI 验证结果（表~\ref{tab:ch4_cpu_tci_results}）、Platform\_TCI 验证结果（表~\ref{tab:ch4_plat_tci_results}），并在最后进行交叉分析。

\textbf{CPU 侧验证结果。}表~\ref{tab:ch4_cpu_tci_results}汇总了不同处理器在各合约层级下的 CPU\_TCI 验证结果。NoFwd 配置在 C1 合约下通过验证（Proven，652 秒），但在 C2 合约下失败（CEX，32.91 秒），说明 NoFwd 防御仅在理想内存平台上安全，无法抵御缓存引入的地址依赖时序差异。Delay 配置在 C1 和 C2 合约下均通过验证，说明 Delay 防御能够在带缓存的平台上保持安全性。Sodor 顺序处理器在 C1 和 C2 合约下均通过验证，且验证时间极短，体现了顺序处理器微架构简单、状态空间小的特征。在 C3 合约下，未启用 PMP 约束的 Sodor 验证失败（CEX，0.44 秒，7 个时钟周期），说明中断控制器引入了数据依赖的时序通道；启用 PMP 约束后验证通过（Proven，5.08 秒），说明 PMP 机制有效阻止了秘密数据对中断行为的影响。

\begin{table}[htb]
    \centering
    \caption{CPU\_TCI 验证结果汇总。PASS 表示验证通过（Proven），FAIL 表示发现反例（CEX），括号内为验证时间。}
    \label{tab:ch4_cpu_tci_results}
    \begin{tabular}{llll}
        \hline
        \textbf{处理器} & \textbf{CPU\_TCI($C_1$)} & \textbf{CPU\_TCI($C_2$)} & \textbf{CPU\_TCI($C_3$)} \\
        \hline
        SimpleOoO NoFwd & PASS (652s) & FAIL (33s) & — \\
        SimpleOoO Delay & PASS (207s) & PASS (8118s) & — \\
        Sodor & PASS (4.82s) & PASS (7.76s) & FAIL (0.44s) \\
        Sodor（PMP） & — & — & PASS (5.08s) \\
        \hline
    \end{tabular}
\end{table}

\textbf{平台侧验证结果。}表~\ref{tab:ch4_plat_tci_results}汇总了不同平台在各合约层级下的 Platform\_TCI 验证结果。常规缓存在 C1 合约下失败（CEX，0.08 秒），因为缓存的命中与缺失导致地址依赖的时序差异，超出 C1 的允许范围；在 C2 合约下通过验证（Proven），因为 C2 允许地址依赖的时序差异。理想内存在 C1 合约下通过验证（Proven），因为其固定延迟设计不产生地址依赖的时序差异；理想内存满足 C1（最强合约），根据格理论的单调性自动满足 C2 及所有更弱的合约层级。中断控制器在 C2 合约下失败（CEX，0.08 秒），因为写入数据影响中断行为超出了 C2 的允许范围；在 C3 合约下通过验证（Proven），因为 C3 允许数据依赖的中断行为。

\begin{table}[htb]
    \centering
    \caption{Platform\_TCI 验证结果汇总。PASS 表示平台满足合约约束（Proven），FAIL 表示平台违反合约（CEX），括号内为验证时间。``蕴含''表示根据格理论单调性，满足更强合约则自动满足更弱合约。}
    \label{tab:ch4_plat_tci_results}
    \begin{tabular}{llll}
        \hline
        \textbf{平台} & \textbf{Platform\_TCI($C_1$)} & \textbf{Platform\_TCI($C_2$)} & \textbf{Platform\_TCI($C_3$)} \\
        \hline
        常规缓存 & FAIL (0.08s) & PASS (0s) & PASS（蕴含） \\
        理想内存 & PASS (0s) & PASS（蕴含） & PASS（蕴含） \\
        中断控制器 & FAIL (0.08s) & FAIL (0.08s) & PASS (0s) \\
        \hline
    \end{tabular}
\end{table}

\textbf{交叉分析。}第\ref{subsec:ch4_decouple}节的不兼容分析预测：若 CPU 侧仅满足更强的合约而平台侧仅满足更弱的合约，则二者不兼容。上述实验数据从正反两方面验证了这一预测。不兼容示例：NoFwd 仅通过 C1 验证（CPU\_TCI($C_1$) PASS，CPU\_TCI($C_2$) FAIL），常规缓存仅通过 C2 验证（Platform\_TCI($C_1$) FAIL，Platform\_TCI($C_2$) PASS），合约需求不匹配，集成后不安全，该定位结果与组合验证的结果（CEX）一致。兼容示例：Delay 通过 C2 验证（CPU\_TCI($C_2$) PASS），常规缓存同样通过 C2 验证（Platform\_TCI($C_2$) PASS），合约层级匹配，根据组合安全性定理集成后安全。此外，C3 合约下同样可以进行交叉定位：Sodor（未启用 PMP）通过 C2 但未通过 C3 验证，中断控制器通过 C3 但未通过 C2 验证，二者在 C3 层级不兼容；启用 PMP 后的 Sodor-S 通过 C3 验证，与中断控制器在 C3 层级兼容。

上述实验结果还揭示了两个重要发现。其一，合约层级能够正确刻画不同平台的时序行为特征：常规缓存满足 C2 但不满足 C1（存在地址依赖时序），理想内存满足 C1（无地址依赖时序），中断控制器满足 C3 但不满足 C2（存在数据依赖中断行为），三类平台被合约层级准确分类。其二，合约层级同样能够正确刻画不同防御机制的强度：NoFwd 仅满足 C1（只能应对理想内存），Delay 满足 C2（能应对缓存平台），Sodor 满足 C2 但不满足 C3（能应对缓存但不能应对中断），Sodor-S（PMP）满足 C3（能应对中断平台），防御机制的安全边界被合约层级准确量化。

\textbf{验证开销分析。}上述实验数据还揭示了验证时间与合约层级、处理器复杂度之间的规律性关系，为工程实践中的验证资源分配提供参考。

从合约层级维度看，同一处理器在更弱的合约下验证时间更长。例如，SimpleOoO Delay 配置在 $C_1$ 合约下验证耗时 207 秒，在 $C_2$ 合约下耗时 8118 秒，增长约 39 倍。这是因为更弱的合约允许更多的平台时序差异（$C_2$ 的 PTCI 引入了额外的无约束响应信号），验证引擎需要探索更大的状态空间以排除所有可能的泄露路径。从处理器复杂度维度看，顺序处理器 Sodor 在 $C_1$ 和 $C_2$ 合约下的验证时间分别为 4.82 秒和 7.76 秒，远低于乱序处理器 SimpleOoO 的对应时间（652 秒和 8118 秒），差距达两个数量级。这体现了顺序处理器微架构简单、不存在推测执行路径的特征，验证引擎需要探索的状态空间显著缩小。

平台侧的验证开销同样值得关注。从表~\ref{tab:ch4_plat_tci_results}可以观察到，所有 Platform\_TCI 验证均在 1 秒以内完成：常规缓存的 $C_1$ 验证仅需 0.08 秒，中断控制器的 $C_2$ 和 $C_3$ 验证分别为 0.08 秒和接近 0 秒。这是因为平台模块（如单条目缓存、中断控制器）的电路规模远小于处理器核心，状态空间有限，且 Platform\_TCI 仅需验证平台自身的时序行为是否符合合约规范，无需同时推理处理器核心的执行逻辑。该结果表明，解耦验证中的平台侧验证不构成验证性能的瓶颈，验证开销几乎全部集中在处理器侧。

% =============================================================================
% 4.4.4 解耦修复验证
% =============================================================================
\subsection{解耦修复验证}
\label{subsec:ch4_eval_fix}

本节验证解耦验证框架能否指导安全修复。基于第\ref{subsec:ch4_eval_locate}节的定位结果，本节设计四种修复方案——修复处理器侧（方案 A）、修复平台侧（方案 B）、两侧同时修复（方案 C）以及 $C_3$ 合约下的修复（方案 D），并通过组合验证（表~\ref{tab:ch4_combined_results}）进行端到端确认。

\textbf{方案 A：修复处理器侧。}将 SimpleOoO 的防御策略从 NoFwd 升级为 Delay。Delay 配置在 C2 合约下通过 CPU\_TCI 验证（Proven，8118 秒），常规缓存在 C2 合约下通过 Platform\_TCI 验证（Proven）。根据组合安全性定理，Delay + 常规缓存的集成安全性得到保证。值得注意的是，当尝试对 Delay + 常规缓存进行传统的组合验证时，验证引擎在 7 天（10080 分钟）内未能得出确定性结论（Undetermined），体现了组合验证面临的状态空间爆炸问题。而本文的解耦验证方法已通过 CPU\_TCI 与 Platform\_TCI 的独立验证给出了安全性保证，无需等待组合验证完成。

\textbf{方案 B：修复平台侧。}将常规缓存替换为理想内存（固定延迟缓存）。NoFwd 配置在 C1 合约下通过 CPU\_TCI 验证（Proven，652 秒），理想内存在 C1 合约下通过 Platform\_TCI 验证（Proven）。根据组合安全性定理，NoFwd + 理想内存的集成安全性得到保证。组合验证确认该结果：NoFwd + 理想内存 → Proven（207 秒）。

\textbf{方案 C：两侧同时修复。}将防御策略升级为 Delay，同时将平台替换为理想内存。Delay 在 C2 合约下通过 CPU\_TCI 验证，理想内存在 C1 合约下通过 Platform\_TCI 验证（C1 强于 C2，根据单调性同样满足 C2）。组合验证确认该结果：Delay + 理想内存 → Proven（185.80 秒）。

\textbf{方案 D：C3 合约下的修复验证。}针对 Sodor 处理器与中断控制器的不兼容问题（Sodor 未通过 C3 验证，中断控制器未通过 C2 验证），通过为 Sodor 启用 PMP 约束进行修复。启用 PMP 后的 Sodor-S 通过 CPU\_TCI($C_3$) 验证（Proven，5.08 秒），中断控制器通过 Platform\_TCI($C_3$) 验证（Proven），根据组合安全性定理二者集成安全。组合验证实验确认该结果：Sodor-S + 中断控制器 → Proven（1.70 秒）。

表~\ref{tab:ch4_fix_results}汇总了各修复方案的验证结果。

\begin{table}[htb]
    \centering
    \caption{解耦修复验证结果。各方案均通过解耦验证保证集成安全性，方案 A 的组合验证因状态空间过大未能在 7 天内完成，体现了解耦验证的可扩展性优势。}
    \label{tab:ch4_fix_results}
    \begin{tabular}{llllll}
        \hline
        \textbf{方案} & \textbf{处理器} & \textbf{平台} & \textbf{合约层级} & \textbf{解耦验证} & \textbf{组合验证} \\
        \hline
        A（修复 CPU） & Delay & 常规缓存 & C2 & PASS & Undetermined (7天) \\
        B（修复平台） & NoFwd & 理想内存 & C1 & PASS & Proven (207s) \\
        C（两侧修复） & Delay & 理想内存 & C2 & PASS & Proven (185.80s) \\
        D（C3 修复） & Sodor-S（PMP） & 中断控制器 & C3 & PASS & Proven (1.70s) \\
        \hline
    \end{tabular}
\end{table}

上述实验表明，设计者可以根据工程约束灵活选择修复方向——修复处理器侧（升级防御机制）、修复平台侧（替换安全平台组件）或同时修复二者。方案 A 的组合验证结果（7 天超时 Undetermined）与解耦验证已给出安全性保证形成鲜明对比，进一步验证了解耦验证相比组合验证在可扩展性上的优势。

% =============================================================================
% 4.4.5 解耦验证与组合验证的对比
% =============================================================================
\subsection{解耦验证与组合验证的对比}
\label{subsec:ch4_eval_scale}

前述实验分别验证了 TCI 插桩的有效性、合约层级的定位能力以及解耦修复的可行性。本节将解耦验证与传统组合验证进行系统对比，从可扩展性与可复用性两个维度论证解耦验证的工程价值。首先通过图~\ref{fig:ch4_cpu_tci_time}和图~\ref{fig:ch4_decoupled_vs_combined}直观对比验证开销，随后分析验证结果的可复用性，最后以兼容性矩阵总结全部实验。

\textbf{验证开销概览。}图~\ref{fig:ch4_cpu_tci_time}展示了不同处理器在各合约层级下的 CPU\_TCI 验证时间。从中可以观察到两个规律：其一，同一处理器在更弱合约下验证时间更长（如 Delay 从 $C_1$ 的 207 秒增至 $C_2$ 的 8118 秒，增长约 39 倍），因为更弱的合约允许更多平台时序差异，验证引擎需探索更大的状态空间；其二，顺序处理器 Sodor 的验证时间远低于乱序处理器 SimpleOoO（秒级 vs 百/千秒级），体现了微架构复杂度对验证开销的直接影响。平台侧验证开销极低，所有 Platform\_TCI 验证均在 1 秒以内完成（常规缓存 0.08 秒，中断控制器 0.08 秒），平台模块的电路规模远小于处理器核心。TCI 插桩本身也不构成额外开销——对比隔离验证（212 秒）与 TCI 插桩验证（32.91 秒），后者反而因使用 Core 模块（去除内部内存）而缩小了验证单元规模。

\begin{figure}[htb!]
    \centering
    \includegraphics[width=0.85\textwidth]{figures/content/4.4 CPU_TCI验证时间对比.pdf}
    \caption{不同处理器在各合约层级下的 CPU\_TCI 验证时间（对数坐标）。更弱的合约允许更多平台时序差异，验证引擎需探索更大的状态空间，验证时间显著增长。}
    \label{fig:ch4_cpu_tci_time}
\end{figure}

\textbf{可扩展性对比。}图~\ref{fig:ch4_decoupled_vs_combined}将解耦验证与组合验证的时间进行直接对比。解耦验证的核心优势在于将联合验证分解为两个独立子问题，规避了组合验证中状态空间随电路规模指数增长的难题。第~\ref{subsec:ch4_eval_fix}节的方案 A 最为直观地体现了这一优势：Delay + 常规缓存的组合验证在 7 天内未能得出确定性结论（Undetermined），而对应的解耦验证通过 CPU\_TCI($C_2$)（8118 秒）与 Platform\_TCI($C_2$)（接近 0 秒）的独立验证已给出安全性保证，总耗时不超过 2.3 小时。组合验证需要同时推理 Delay 处理器核心的流水线状态与常规缓存的缓存状态，二者的状态空间乘积导致验证引擎无法在合理时间内收敛；解耦验证则将同一问题拆分为两个规模可控的子问题，状态空间远小于联合状态空间。

\begin{figure}[htb!]
    \centering
    \includegraphics[width=0.9\textwidth]{figures/content/4.4 解耦验证vs组合验证时间对比.pdf}
    \caption{解耦验证与组合验证的时间对比（对数坐标）。Delay + 常规缓存场景下组合验证 7 天超时，而解耦验证在 2.3 小时内完成，体现了显著的可扩展性优势。}
    \label{fig:ch4_decoupled_vs_combined}
\end{figure}

对于小规模设计，解耦验证在绝对时间上不一定优于组合验证。例如，Sodor-S + 中断控制器的组合验证耗时 1.70 秒，而对应的解耦验证总耗时 5.08 秒。然而，随着处理器复杂度的增长（如从 Sodor 到 SimpleOoO），组合验证的时间开销急剧膨胀（从秒级到 7 天超时），而解耦验证的增长幅度相对可控（从 5 秒到 2.3 小时），其可扩展性优势随设计规模增大而愈发显著。

\textbf{可复用性对比。}解耦验证的另一核心优势在于验证结果可跨设计变更复用。在本文的实验中，2 款处理器（NoFwd、Delay）与 2 种缓存平台（常规缓存、理想内存）共产生 $2 \times 2 = 4$ 种集成组合。传统组合验证需要对每种组合独立进行完整验证（共 4 次，其中 1 次 7 天超时），任何一侧的设计变更都要求重新执行所有涉及该组件的组合验证。而解耦验证仅需 2 次 CPU\_TCI 与 2 次 Platform\_TCI 共 4 次独立验证，即可通过兼容性矩阵覆盖全部 4 种组合的安全性判定。当新增一种平台组件时，组合验证需要新增与每款处理器的组合验证（本实验中为 2 次），而解耦验证仅需执行 1 次新平台的 Platform\_TCI 验证（通常不到 1 秒），已有的 CPU\_TCI 结果可直接复用。类似地，当处理器更换防御机制时，仅需重新执行该处理器的 CPU\_TCI 验证，已有的 Platform\_TCI 结果不受影响。

\textbf{兼容性矩阵。}表~\ref{tab:ch4_compatibility}将全部实验的 CPU 侧与平台侧验证结果交叉汇总为兼容性矩阵。该矩阵的每个单元格均无需执行对应的组合验证，而是直接由 CPU\_TCI 与 Platform\_TCI 的独立结果推导得出，直观展示了解耦验证的可复用性。

\begin{table}[htb]
    \centering
    \caption{处理器核心与平台的兼容性矩阵。$\checkmark$ 表示二者在共同满足的合约层级下兼容（集成安全），$\times$ 表示不兼容（存在时序泄露风险），``—''表示未测试。}
    \label{tab:ch4_compatibility}
    \begin{tabular}{llll}
        \hline
        \textbf{处理器} & \textbf{常规缓存} & \textbf{理想内存} & \textbf{中断控制器} \\
        & （满足 $C_2$） & （满足 $C_1$） & （满足 $C_3$） \\
        \hline
        NoFwd（满足 $C_1$） & $\times$ & $\checkmark$（$C_1$） & — \\
        Delay（满足 $C_2$） & $\checkmark$（$C_2$） & $\checkmark$（$C_1$） & — \\
        Sodor（满足 $C_2$） & $\checkmark$（$C_2$） & $\checkmark$（$C_1$） & $\times$ \\
        Sodor-S / PMP（满足 $C_3$） & $\checkmark$（$C_2$） & $\checkmark$（$C_1$） & $\checkmark$（$C_3$） \\
        \hline
    \end{tabular}
\end{table}

综上，解耦验证在可扩展性与可复用性上相比组合验证具有显著优势。在可扩展性方面，解耦验证将联合状态空间分解为独立的子空间，在 Delay + 常规缓存场景下将组合验证的 7 天超时缩短至 2.3 小时。在可复用性方面，CPU\_TCI 与 Platform\_TCI 结果可独立复用，新增或替换组件时仅需增量验证变更侧，避免了组合验证中``全部重来''的高昂代价。




% ============================================================================
% 4.5 反例分析
% ============================================================================
\section{反例分析}
\label{sec:ch4_cex_analysis}

形式化验证引擎在发现安全属性违反时，会返回一个具体的反例（Counterexample, CEX），包含导致泄露的完整指令序列、秘密数据赋值以及逐周期的信号翻转轨迹。反例不仅表示验证失败，更提供了一个完整的攻击场景描述，可直接用于理解泄露机制、定位现有防御的薄弱环节，并评估其可利用性。本节对第\ref{sec:ch4_experiment}节实验中发现的典型反例进行深入分析，重点剖析攻击链如何绕过现有设计，以及这些违规对构造实际攻击的启示。

表~\ref{tab:ch4_cex_summary}汇总了实验中发现的五组典型反例。

\begin{table}[htb]
    \centering
    \caption{典型反例汇总。每组反例对应一种时序泄露机制，反例周期数表示验证引擎找到泄露所需的最短执行轨迹。}
    \label{tab:ch4_cex_summary}
    \begin{tabular}{lllrl}
        \hline
        \textbf{类别} & \textbf{验证对象} & \textbf{合约} & \textbf{周期数} & \textbf{泄露机制} \\
        \hline
        CPU 侧 & NoFwd CPU\_TCI($C_2$) & $C_2$ & 16 & 推测 load 缓存时序泄露 \\
        CPU 侧 & Sodor CPU\_TCI($C_3$) & $C_3$ & 7 & 秘密数据触发中断差异 \\
        平台侧 & 常规缓存 Platform\_TCI($C_1$) & $C_1$ & 1 & 地址依赖响应时序 \\
        平台侧 & 中断控制器 Platform\_TCI($C_2$) & $C_2$ & 2 & 写入数据触发中断 \\
        组合 & NoFwd + 常规缓存 & — & 16 & 推测 load 经缓存泄露 \\
        \hline
    \end{tabular}
\end{table}

\subsection{CPU 侧反例：推测 load 绕过 NoFwd 防御}
\label{subsec:ch4_cex_nofwd}

NoFwd CPU\_TCI($C_2$) 反例在 16 个时钟周期内构造了一条完整的 Spectre 型攻击链，揭示了 NoFwd 防御策略在面对缓存型平台时的根本性缺陷。该反例的攻击链可分为四个阶段。

\textbf{阶段一：推测路径进入。}流水线执行一条分支指令，分支预测器将执行引导至推测路径。在 SimpleOoO 的 4 级流水线中，分支结果尚未提交，后续指令已被推测发射。

\textbf{阶段二：推测 load 发射。}在推测路径上，一条 load 指令以秘密数据作为地址被发射至内存总线。NoFwd 防御的设计意图是``允许推测 load 发射但不将结果转发给后续指令''，即通过切断数据转发通路来防止秘密值污染架构状态。然而，NoFwd 未阻止推测 load 的地址信号到达处理器核心的内存接口。

\textbf{阶段三：平台时序差异注入。}在 $C_2$ 合约假设下，平台的响应时序依赖于访存地址（如缓存命中 1 周期、缺失 3 周期）。两个处理器副本使用不同的秘密值，产生不同的访存地址，进而收到不同延迟的平台响应。至此，秘密信息已通过地址→响应延迟的路径从处理器核心泄露至时序域。

\textbf{阶段四：时序差异传播与观测。}平台响应的时序差异沿流水线反向传播：不同的响应延迟导致后续指令的流水线占用时间不同，最终体现为非推测指令提交时间的差异。该差异在可观测输出（提交时序）上被攻击者捕获。

\textbf{NoFwd 失效的根因。}NoFwd 的防御逻辑仅关注数据通路——阻止推测 load 的结果被转发。但它忽视了控制通路和时序通路：推测 load 的地址仍然到达内存接口，而平台的响应时序携带了地址依赖的信息，这些时序差异沿流水线回传并影响提交时间。NoFwd 保护的是``数据不泄露''，但未保护``时序不泄露''。

\textbf{攻击启示与防御参考。}该反例与经典的 Spectre-v1（边界检查绕过）攻击具有相同的泄露原语：推测访存以秘密为地址 → 缓存命中/缺失差异 → 时序侧信道。这表明，在带缓存的平台上，仅防御数据转发是不够的。Delay 策略通过在推测阶段完全抑制 load 发射，从源头切断了推测地址到达平台的路径，因此能够在 $C_2$ 合约下通过验证。组合验证实验（NoFwd + 常规缓存 → CEX，16 周期）产生了周期数和泄露机制均一致的反例，进一步验证了 TCI 插桩对真实缓存行为的等效模拟。

\subsection{CPU 侧反例：中断通道绕过顺序执行假设}
\label{subsec:ch4_cex_sodor}

Sodor CPU\_TCI($C_3$) 反例仅需 7 个时钟周期即可触发，揭示了一种不依赖推测执行的全新时序泄露向量。该反例表明，即使在最简单的顺序处理器上，平台交互仍可引入安全性隐患。

\textbf{攻击链。}Sodor 为 2 级流水线顺序处理器，不存在推测执行路径，因此免疫 Spectre 类攻击。然而，反例构造了以下攻击序列：一条已提交的 store 指令将秘密数据写入外设地址范围内的内存映射寄存器（中断控制器），不同的写入数据触发不同的中断响应行为。在 $C_3$ 合约假设下，中断信号差异导致两个处理器副本进入不同的执行路径——一个副本响应中断进入中断处理例程，另一个副本继续正常执行，产生显著的提交时序差异。

\textbf{Sodor 失效的根因。}Sodor 的设计未对内存映射外设的访问施加隔离约束——任何 store 指令均可写入中断控制器的地址范围，包括携带秘密数据的 store。这并非 Sodor 的设计缺陷，而是反映了一个系统层面的盲区：传统处理器安全验证聚焦于推测执行防御，而忽略了非推测路径上平台交互引入的时序通道。

\textbf{攻击启示与防御参考。}该反例揭示了中断时序侧信道的威胁模型：攻击者无需利用推测执行，只需控制一条 store 指令的写入数据和目标地址即可完成泄露。该攻击的触发条件极为简单（7 个周期，无需分支预测或缓存状态准备），且适用于所有具备内存映射中断接口的处理器。PMP（物理内存保护）机制通过限制秘密数据可写入的地址范围，从软件层面阻断了``秘密数据 → 外设地址''的信息流，有效消除了该攻击路径。该结论对嵌入式系统的安全设计具有直接参考价值：即使处理器核心本身经过严格验证，未受保护的外设访问仍可构成泄露通道。

\subsection{平台侧反例}
\label{subsec:ch4_cex_platform}

平台侧反例刻画了平台组件自身的时序行为与合约层级的对应关系，帮助设计者理解每类平台的时序特征边界。

\textbf{常规缓存 Platform\_TCI($C_1$) 反例（1 周期）。}该反例是所有反例中最简单的，仅需 1 个周期即可触发。两个缓存副本接收不同的访存地址，命中副本在 1 周期内返回响应，缺失副本需要 3 周期，响应时序产生差异。在 $C_1$ 合约下地址不应影响响应时序，因此该差异构成违规。该反例直观地表明常规缓存不满足 $C_1$ 合约，其地址依赖的命中/缺失行为本质上就是缓存时序侧信道的硬件根因，处理器设计必须在缓存平台假设（$C_2$）下进行安全性验证。

\textbf{中断控制器 Platform\_TCI($C_2$) 反例（2 周期）。}两个中断控制器副本接收相同的地址但不同的写入数据，一个副本的写入数据触发了中断信号跳变，另一个未触发，产生中断信号的差异。在 $C_2$ 合约下写入数据不应影响响应时序和中断信号，因此该差异构成违规。该反例表明，中断控制器的``数据 → 中断''信息流超出了 $C_2$ 的允许范围，需要 $C_3$ 合约才能正确描述。对攻击者而言，该信息流意味着通过控制写入数据即可操纵中断行为，进而影响处理器的执行时序。

\subsection{反例分析小结}

上述反例分析表明，时序合约框架不仅能够判定验证通过与否，还能通过反例提供具体的攻击场景描述，帮助设计者理解泄露的触发条件和传播路径，评估攻击的可利用性。CPU 侧反例揭示了现有防御机制的盲区：NoFwd 反例表明仅防御数据转发不足以抵御缓存时序侧信道，Sodor 反例表明顺序执行不能豁免外设交互引入的时序泄露。平台侧反例则刻画了平台组件的时序行为边界，为合约层级的选择提供了直接依据。这些反例不仅服务于当前设计的安全评估，还可作为未来防御策略设计的参考基准——新的防御方案应确保能够阻断上述反例中识别出的泄露路径。


% ============================================================================
% 4.6 本章小结
% ============================================================================
\section{本章小结}
\label{sec:ch4_summary}

本章针对现有形式化验证方案在处理器核心隔离验证中存在的安全性盲区以及组合验证面临的可扩展性瓶颈，提出了基于片上系统时序合约的解耦验证框架。

在时序合约方面，本章通过分析理想内存、缓存、内存映射外设三类常见平台的时序信号特征，提出了片上系统时序合约的形式化定义方法，分别定义了 C1、C2、C3 三种合约层级，系统化描述了处理器核心与平台之间可能存在的时序泄露关系，并基于格理论建立了合约之间的强弱偏序关系，为分层验证提供了理论依据。

在验证方法方面，本章设计了时序合约插桩方法，通过 XOR 差异检测、sticky-one 算子与插桩 MUX 控制三个核心组件，将合约定义编码为可综合的电路逻辑。在此基础上，本章构建了双向解耦验证策略，将时序合约插桩分解为处理器侧（CPU\_TCI）与平台侧（Platform\_TCI）两类，分别对处理器核心和平台组件进行独立验证，通过组合安全性定理保证集成后的安全性。

实验结果表明，本章提出的解耦验证框架在多款 RISC-V 处理器（SimpleOoO、Sodor）和多种平台组件（常规缓存、理想内存、中断控制器）上取得了有效验证，具体体现在以下四个方面：其一，覆盖范围从处理器核心拓展至处理器与平台的交互，时序合约插桩在无需平台 RTL 代码的条件下检出了隔离验证无法发现的时序泄露；其二，时序合约插桩检出的泄露与组合验证一致，且验证速度提升约 1.8 倍，验证过程无需获取平台 RTL 代码；其三，在组合验证 7 天超时的场景下，解耦验证仍能在 2.3 小时内给出安全性保证，体现了显著的可扩展性优势；其四，验证结果可跨设计变更复用，新增平台组件仅需执行一次 Platform\_TCI 验证，已有的 CPU\_TCI 结果可直接复用，验证成本与组合数量解耦。
